\chapter{Contact dermatitis}

\section*{Definition and Overview}
Contact dermatitis is an inflammatory skin condition characterised by an erythematous, pruritic, and sometimes vesicular eruption resulting from direct skin contact with an external substance. It is broadly categorised into two main forms: irritant contact dermatitis (ICD) and allergic contact dermatitis (ACD). Irritant contact dermatitis is a non-immunological reaction caused by direct chemical or physical damage to the epidermal barrier. In contrast, allergic contact dermatitis is a type IV delayed-type hypersensitivity reaction mediated by T lymphocytes in individuals who have been previously sensitised to a specific allergen. Distinguishing between these two forms is essential for formulating effective clinical management plans and implementing appropriate preventive measures.

\section*{Epidemiology}
Contact dermatitis is highly prevalent and represents one of the most common occupational diseases globally. It is estimated that approximately 15\% to 20\% of the general population will experience contact dermatitis at some point in their lives. Irritant contact dermatitis is the more common subtype, accounting for approximately 80\% of all contact dermatitis cases. It frequently affects individuals in occupations requiring frequent wet work or exposure to chemicals, such as healthcare workers, hairdressers, cleaners, food handlers, and construction workers. Allergic contact dermatitis can occur at any age and is commonly triggered by nickel, fragrances, cosmetic preservatives, and topical medications. In Australia, contact dermatitis accounts for a substantial proportion of work-related skin conditions, causing significant morbidity and economic burden due to lost productivity.

\section*{Aetiology and Pathophysiology}
The underlying mechanisms of contact dermatitis differ based on whether the reaction is triggered by physical or chemical irritation or by an allergen-induced immune response.

\begin{itemize}
    \item \textbf{Chemical irritants:} Direct cellular damage to keratinocytes and disruption of the intercellular lipid barrier by harsh chemicals, such as soaps, detergents, solvents, acids, and alkalis.
    \item \textbf{Physical and environmental factors:} Micro-trauma caused by low humidity, cold temperatures, wind, friction, and prolonged exposure to water (wet work), which compromises the skin's protective barrier.
    \item \textbf{Allergen-induced sensitisation:} Initial exposure to small molecules called haptens (such as nickel, chromate, fragrances, or preservatives like methylisothiazolinone) that bind to epidermal proteins to form immunogenic complexes.
    \item \textbf{Type IV hypersensitivity reaction:} Subsequent exposure to the allergen triggers antigen presentation by Langerhans cells to sensitised memory T lymphocytes, leading to a localised inflammatory cascade 24 to 72 hours later.
    \item \textbf{Genetic barrier susceptibility:} Pre-existing epidermal barrier defects, such as those caused by filaggrin gene mutations or a history of atopic dermatitis, which increase vulnerability to both irritant and allergic contact dermatitis.
\end{itemize}

\section*{Clinical Features}
The clinical manifestations of contact dermatitis primarily occur at the site of contact, although severe allergic reactions may spread beyond the initial area of exposure.

\begin{itemize}
    \item \textbf{Pruritus and burning:} Intense itching, which is highly characteristic of allergic contact dermatitis, or a prominent burning and stinging sensation, which is more commonly associated with irritant contact dermatitis.
    \item \textbf{Erythema and oedema:} Redness and swelling of the skin, often presenting in a shape or pattern that reflects the offending object or substance (e.g., a linear rash from a plant brush or a circular rash under a watch).
    \item \textbf{Vesiculation and bullae:} The formation of small blisters (vesicles) or large blisters (bullae) that may rupture, ooze, and crust in acute cases.
    \item \textbf{Scaling and fissuring:} Dryness, scaling, and painful cracks (fissures) in the skin, which are characteristic of chronic contact dermatitis, particularly on the hands.
    \item \textbf{Lichenification:} Thickened, leathery skin with exaggerated skin lines, resulting from chronic rubbing and scratching of pruritic areas.
\end{itemize}

\section*{Differential Diagnosis}
Contact dermatitis must be differentiated from other common eczematous and inflammatory skin disorders.

\begin{itemize}
    \item \textbf{Atopic dermatitis:} An endogenous, chronic pruritic dermatitis typically presenting in early childhood with a characteristic flexural distribution and a personal or family history of atopic diseases.
    \item \textbf{Seborrheic dermatitis:} An inflammatory condition characterised by greasy, yellowish scales concentrated in sebum-rich areas, such as the scalp, nasolabial folds, and chest.
    \item \textbf{Dyshidrotic eczema (pompholyx):} A vesicular hand and foot dermatitis characterised by deep-seated, highly pruritic blisters on the palms, soles, and lateral aspects of the digits.
    \item \textbf{Psoriasis vulgaris:} Well-demarcated plaques covered with silvery scales, typically located on extensor surfaces (elbows and knees) and often associated with nail pitting.
    \item \textbf{Tinea infections:} Superficial fungal infections (such as tinea pedis or tinea manuum) that present with annular plaques, scaling borders, and can be confirmed with skin scrapings.
\end{itemize}

\section*{When to Seek Medical Care}
Patients should seek medical evaluation from a general practitioner or dermatologist if the skin rash is severe, painful, rapidly spreading, fails to improve with over-the-counter hydrocortisone creams, or if signs of secondary bacterial infection (such as pus, warmth, or red streaks) develop.

\textbf{Call 000 (or local emergency services) for:}
\begin{itemize}
    \item Swelling of the face, lips, tongue, or throat that may compromise the airway.
    \item Difficulty breathing, wheezing, or severe shortness of breath.
    \item A sudden, widespread rash accompanied by high fever, chills, or systemic illness.
    \item Extensive blistering or painful peeling of the skin covering a large portion of the body.
\end{itemize}

\section*{Diagnosis and Investigations}
Diagnosis is primarily based on a thorough clinical history, physical examination, and, when indicated, specialised dermatological investigations to identify specific allergens.

\begin{itemize}
    \item \textbf{Clinical history and occupational review:} Detailed questioning regarding daily activities, hobbies, cosmetic products, work environment, and the relationship of the rash to periods of absence from work.
    \item \textbf{Patch testing:} The gold standard investigation for allergic contact dermatitis, involving the application of standardised allergen panels (e.g., European Standard Series) to the back for 48 hours, with subsequent readings at 48 and 72 to 96 hours.
    \item \textbf{Repeated open application test (ROAT):} A simple test where the patient applies a suspected product (such as a cosmetic or lotion) to a small area of normal skin twice daily for up to one week to monitor for a localised reaction.
    \item \textbf{Skin scrapings for microscopy:} Microscopic examination of skin scrapings prepared with potassium hydroxide (KOH) to rule out superficial fungal infections (tinea).
    \item \textbf{Skin biopsy:} Occasionally performed in atypical, severe, or treatment-resistant cases to rule out other dermatological conditions, such as cutaneous T-cell lymphoma or psoriasis.
\end{itemize}

\section*{Management}
The primary objective of management is the identification and complete avoidance of the offending substance, coupled with pharmacological treatments to suppress cutaneous inflammation and restore barrier function.

\begin{itemize}
    \item \textbf{Avoidance and substitution:} Eliminating contact with the identified allergen or irritant and substituting it with hypoallergenic, fragrance-free alternatives.
    \item \textbf{Topical corticosteroids:} The mainstay of pharmacological therapy to reduce inflammation and itching; the potency of the corticosteroid is selected based on the severity of the rash and the skin thickness of the affected area.
    \item \textbf{Emollients and barrier creams:} Regular and liberal application of thick, fragrance-free moisturising creams or ointments to repair the compromised epidermal barrier and prevent water loss.
    \item \textbf{Topical calcineurin inhibitors:} Non-steroidal anti-inflammatory agents (such as tacrolimus ointment or pimecrolimus cream) that are particularly useful for sensitive skin sites like the face and eyelids to avoid corticosteroid-induced skin atrophy.
    \item \textbf{Oral corticosteroids:} A short course of systemic corticosteroids (e.g., prednisolone) for patients with severe, widespread contact dermatitis or significant facial involvement.
    \item \textbf{Phototherapy and systemic immunosuppressants:} Narrowband UVB phototherapy or oral immunosuppressive agents (e.g., methotrexate, ciclosporin) for chronic, refractory cases that fail to respond to standard topical therapies.
\end{itemize}

\section*{Complications}
Failure to manage contact dermatitis or identify the offending agent can lead to several complications.

\begin{itemize}
    \item \textbf{Secondary bacterial infection:} Impetiginisation or cellulitis, commonly caused by \textit{Staphylococcus aureus} or \textit{Streptococcus pyogenes} entering through compromised, scratched skin.
    \item \textbf{Chronic hand eczema:} A long-term, debilitating condition characterised by persistent skin thickening, scaling, and painful fissuring that limits hand function.
    \item \textbf{Post-inflammatory dyschromia:} Wavelength-dependent hyperpigmentation or hypopigmentation at the site of inflammation, which can persist for months after the rash resolves.
    \item \textbf{Occupational and lifestyle disruption:} Severe distress, loss of employment, or necessity of career change due to inability to perform occupational tasks without triggering the condition.
    \item \textbf{Sleep disturbance and psychological distress:} Persistent, intense pruritus leading to insomnia, fatigue, anxiety, and depression.
\end{itemize}

\section*{Prognosis and Outcomes}
The prognosis for contact dermatitis is generally excellent if the offending irritant or allergen is successfully identified and avoided. Acute contact dermatitis typically resolves within two to three weeks with appropriate topical corticosteroid therapy and barrier repair. However, the prognosis is less favourable for patients with occupational contact dermatitis where complete avoidance of triggers is challenging, or in individuals with underlying atopic dermatitis. In these cases, the disease can become chronic and recurrent, requiring persistent self-care, strict barrier protection, and intermittent use of topical anti-inflammatory medications.

\section*{Prevention and Self-Care}
Preventative strategies focus on protecting the skin barrier and minimising exposure to known chemical and physical irritants.

\begin{itemize}
    \item \textbf{Use of personal protective equipment:} Wearing appropriate protective gloves (such as nitrile gloves for chemical handling and wet work) and long sleeves to prevent skin exposure.
    \item \textbf{Application of barrier creams:} Applying specialized barrier creams prior to starting work activities to help shield the skin from mild irritants.
    \item \textbf{Gentle skin cleansing:} Washing the skin immediately with a mild, soap-free cleanser and lukewarm water following accidental exposure to a known irritant or allergen.
    \item \textbf{Liberal moisturisation:} Applying fragrance-free emollients frequently throughout the day, particularly immediately after washing the hands, to support the skin's lipid barrier.
    \item \textbf{Diligent ingredient review:} Checking ingredient lists on cosmetics, toiletries, and household products to identify and avoid known personal allergens.
\end{itemize}

\section*{Sources and Further Reading}
The following reputable organisations provide further information and clinical resources on contact dermatitis:

\begin{itemize}
    \item Australasian College of Dermatologists. \textit{Patient resources, fact sheets, and information on contact allergies and patch testing.} https://www.dermcoll.edu.au/
    \item British Association of Dermatologists. \textit{Patient information leaflets on contact dermatitis, occupational skin disease, and patch testing.} https://www.bad.org.uk/
    \item DermNet. \textit{Comprehensive clinical articles, image libraries, and management guidelines for irritant and allergic contact dermatitis.} https://dermnetnz.org/
    \item Mayo Clinic. \textit{Overview of contact dermatitis, symptoms, causes, and home care remedies.} https://www.mayoclinic.org/
    \item NHS. \textit{General health advice, prevention strategies, and treatment options for contact dermatitis.} https://www.nhs.uk/
\end{itemize}
